\section{Introduction}
Exact input transformations provide a controlled way to study what a neural representation preserves. A task answer may transform predictably when an input is inverted, complemented, translated, or differentiated. Yet matching transformed answers, aligning hidden vectors, and predicting an unseen composition are different achievements. A representation can support an accurate first-step answer while omitting information required by the next operation.

We ask whether \emph{the correctness of a supplied state correspondence} improves held-out composition when output supervision is held fixed. The distinction matters because geometric supervision ordinarily changes both the state target and the answer information available during training. We use exact finite mathematical worlds to construct incorrect correspondences that preserve visible gold answers, endpoint marginals, input exposure, and correct teacher-output supervision. The intervention changes which intermediate hidden state is treated as the geometric target.

Our primary evidence comes from a final confirmation with five new source initializations in a finite matrix world. Encoders first receive ordinary task training and correct single-generator preparation. We then freeze each encoder, its second operator, and its readout, and cross first-operator activation with correct versus matched-incorrect correspondence. This is a conditional intervention on relation-prepared representations. It does not test whether ordinary training discovers algebra spontaneously.

Three findings form the paper's argument. First, under a fixed budget, correctly paired GELU operators support more accurate unseen composition than their matched-incorrect counterparts. Second, predicted components invisible to the fixed first-step linear readout can transmit this benefit: swapping those components improves a linear recipient while preserving every first-step logit. Third, adequately fitting the same affine function family under the same single-step objective recovers most of the functional correctness benefit. Consequently, the original budget-controlled nonlinear advantage does not establish an additional expressivity advantage.

Permutation studies provide a second instance of correspondence-sensitive composition and motivate the intermediate-state diagnostics. Initial joint affine matrix experiments, an irreversible polynomial experiment, and a null-only-loss intervention establish important boundaries. Their architectures, supervision directions, source counts, and fitting budgets differ; they are supporting studies, rather than a pooled cross-domain confirmation of one universal mechanism.

The contribution is a controlled empirical account of \emph{when geometric correspondence is functionally useful and how that conclusion can be checked}. It combines answer-matched correspondence interventions, predicted-component transfer, independent-source confirmation, and same-family convergence diagnostics. Readers interested in compositional representation learning can use this combination to distinguish answer preservation, downstream state sufficiency, and finite-budget optimization. We make no claim that each diagnostic alone is new, or that the resulting mechanism holds across all mathematical systems.
